% ===================== INTRODUCCION =====================
\section{Introducción}
\label{sec:intro}

Este documento presenta la evaluación y validación del diseño hidráulico y de riesgos operativos elaborado para la reubicación por Perforación Horizontal Dirigida (PHD) del salmueroducto e infraestructura de condensados de BRINSA S.A. en cercanías del Humedal Arrieros y el parque Jaime Duque, en la coordenada Km 16.5 de la red.

% ─────────────────────────────────────────────────────────────────────────────
\subsection{Antecedentes}
\label{ssec:intro_antecedentes}

El salmueroducto e infraestructura de condensados transporta fluidos esenciales para la producción de sal refinada y cloro-álcali entre la mina de Zipaquirá y la planta industrial en Tocancipá. Debido a consideraciones ecológicas y de conservación de la biodiversidad, la autoridad ambiental del departamento requirió un desvío y reubicación del trazado de las tuberías superficiales en el tramo Km 16.5 para alejar las líneas de la zona inundable del Humedal Arrieros. Para responder a este requerimiento regulatorio protegiendo la superficie, el equipo de Mantenimiento de BRINSA S.A. propuso una alternativa civil basada en Perforación Horizontal Dirigida (PHD) para enterrar y proteger el tramo en una profundidad de \SI{15.60}{\meter}.

% ─────────────────────────────────────────────────────────────────────────────
\subsection{Descripción del Sistema}
\label{ssec:intro_descripcion}

El sistema bajo estudio se compone de dos líneas de conducción paralelas:
1. Una línea de salmuera saturada al \SI{26}{\percent} p/p de \text{NaCl} a una temperatura de \SI{15}{\celsius}, con caudal volumétrico de \SI{300.00}{\cubic\meter\per\hour} y diámetro nominal de \SI{12}{in}.
2. Una línea de condensado de vapor industrial a una temperatura de \SI{20}{\celsius}, con caudal volumétrico de \SI{300.00}{\cubic\meter\per\hour} y diámetro nominal de \SI{12}{in}.
El material base planteado para ambas tuberías es Polietileno de Alta Densidad (HDPE) PE100 SDR 13.6 (condensado) y SDR 11 (salmuera).

% ─────────────────────────────────────────────────────────────────────────────
\subsection{Motivación del Estudio}
\label{ssec:intro_motivacion}

La profundización vertical de \SI{15.60}{\meter} a lo largo del perfil de cruce de \SI{450.00}{\meter} crea una configuración geométrica de sifón invertido. La Dirección de Ingeniería y Proyectos de BRINSA S.A. ha solicitado a DML S.A.S. una validación hidráulica preliminar para determinar:
1. Si la mayor longitud física inclinada de la tubería (\SI{452.26}{\meter}) y las curvas de la perforación dirigida generan pérdidas de carga dinámicas que afecten la succión de la estación Km8 y Km17.
2. Cuáles son las sobrepresiones mecánicas inducidas por la columna hidrostática vertical en el fondo de la perforación y qué riesgos operativos de acumulación de sólidos o entrampamiento de aire se derivan del sifón invertido.

% ─────────────────────────────────────────────────────────────────────────────
\subsection{Estructura del Documento}
\label{ssec:intro_estructura}

El informe técnico se estructura de la siguiente manera: la Sección 2 presenta las Bases de Diseño que rigen los fluidos, dimensiones y la Sabana de Bogotá; la Sección 3 detalla la Metodología y ecuaciones hidráulicas implementadas; la Sección 4 expone los Resultados numéricos de velocidad y pérdidas de presión dinámicas y estáticas; la Sección 5 contiene el Análisis técnico crítico de los riesgos operativos y mecánicos identificados; finalmente, las Secciones 6 y 7 presentan las Conclusiones y Recomendaciones de ingeniería.
